\section{Ray-marching - Maximum iterations versus distance threshold condition}\label{ray-marching---maximum-number-of-iterations-versus-distance-threshold-condition}

The \emph{ray marching distance threshold}\index{ray marching!distance threshold} is the condition where a ray point
marching along the ray comes within a specified distance from the fractal
surface and the ray-marching stops. This controls the size of the detail in the
image, and is normally set to vary such that greater detail is obtained for the
surface closest to the camera, while in the further regions of the fractal, the
distance threshold will be larger such that only bigger details are visible.
Enabling \emph{Constant Detail Size}\index{ray marching!constant detail size} on the \emph{Rendering Engine} tab will
make the distance threshold uniform across the entire image.

For more information on Distance Estimation and its relationship to the distance threshold, see chapter \ref{distance-estimation}.

There are two modes of controlling the ray-marching termination for each image pixel.

\subsection{Mode 1: Stop at distance threshold (default)}\label{mode-1-stop-at-distance-threshold}

When \emph{Stop at maximum iteration} is \textbf{disabled} (default), ray-marching stops when the estimated distance to the fractal surface falls below the distance threshold. The distance threshold is calculated dynamically based on the distance from the camera and the \emph{Detail Level} setting. This is the most efficient mode because the fractal iteration loop typically stops on the \emph{Bailout}\index{termination condition!bailout} condition before reaching \emph{formula\_maxiter} (see page \pageref{bailout-maxiter}), making rendering much faster.

\subsection{Mode 2: Stop at maximum iterations}\label{mode-2-stop-at-maximum-iterations}

When \emph{Stop at maximum iteration} is \textbf{enabled}, ray-marching no longer stops at the distance threshold. Instead, the detail level is controlled by the \emph{Max. fractal iterations} (\emph{formula\_maxiter}) parameter. In this mode, the distance threshold is proportional to the distance from the camera, and the fractal iteration loop runs to the full \emph{formula\_maxiter} count on the fractal surface (when bailout is not reached).

\textbf{Important note:} \emph{Stop at maximum iteration} does not control the fractal iteration loop directly. It controls the ray-marching termination condition. The fractal iteration loop always runs to achieve \emph{Bailout}\index{termination condition!bailout}, then if bailout is not reached the iteration stops at \emph{formula\_maxiter} (see page \pageref{bailout-maxiter}).

\clearpage

\subsection{Comparison of the two modes}

\twoImagesWithTwoCaptionsFullWidth{img/manual/media/stop_raymarching_at_disttrhersh.png}
{Mode 1: Stop ray-marching at distance threshold}
{stop_raymarching_at_disttrhersh}
{img/manual/media/stop_raymarching_at_maxiter.png}
{Mode 2: Stop ray-marching at Maxiter}
{stop_raymarching_at_maxiter}{H}

In \textbf{Mode 1} (left figure), ray-marching stops at the distance threshold. In most cases the fractal iteration loop stops on the bailout condition because away from the surface it is not possible to reach \emph{formula\_maxiter}. This makes rendering of fractals much faster.

In \textbf{Mode 2} (right figure), ray-marching continues until the maximum number of ray-marching steps is reached (the ray-marching distance threshold is ignored). On the fractal surface, the maximum number of fractal iterations is calculated (when bailout is not reached).

\subsection{Low \emph{formula\_maxiter} behavior}

Figure \ref{stop-raymarching-at-bailout-low-maxiter} shows what happens when \emph{formula\_maxiter} is set to a low value (e.g., 4). Even if \emph{formula\_maxiter}\index{termination condition!maxiter} is reached, the ray-marching continues until the ray-marching distance threshold is reached (in Mode 1).

\twoImagesWithTwoCaptionsFullWidth{img/manual/media/stop_raymarching_at_disttrhersh_iter4.png}
{Mode 1: Stop ray-marching at bailout with low Maxiter}
{stop-raymarching-at-bailout-low-maxiter}
{img/manual/media/stop_raymarching_at_maxiter_iter4.png}
{Mode 2: Stop ray-marching at maxiter with low Maxiter}
{stop-raymarching-at-maxiter-low-maxiter}{h}

\clearpage

\subsection{Constant Detail Size}

Enabling \emph{Constant Detail Size}\index{ray marching!constant detail size} on the \emph{Rendering Engine} tab will make the distance threshold uniform across the entire image. This takes longer to render, but gives a more accurate representation of detail at different distances, and can be used to address color variation over distance (figure \ref{constant_detail_size}).

\simpleImageWithCaptionHalfWidth{img/manual/media/constant_detail_size.jpg}
{Constant detail size}
{constant_detail_size}{H}
